\documentclass[12pt,a4paper]{article}
\usepackage[utf8]{inputenc}
\usepackage[T1]{fontenc}
\usepackage{geometry}
\usepackage{hyperref}
\usepackage{color}
\usepackage{listings}
\usepackage{longtable}
\usepackage{booktabs}
\usepackage{array}
\usepackage{enumitem}
\usepackage{fancyhdr}
\usepackage{titlesec}
\usepackage{xcolor}

\geometry{margin=2.5cm}

% Color definitions
\definecolor{linkblue}{RGB}{0,0,255}
\definecolor{citegreen}{RGB}{0,128,0}
\definecolor{warningred}{RGB}{255,0,0}
\definecolor{noteorange}{RGB}{255,165,0}

% Hyperref setup
\hypersetup{
    colorlinks=true,
    linkcolor=linkblue,
    citecolor=citegreen,
    urlcolor=linkblue,
    pdftitle={CryptoGL Official Specifications and Test Vectors},
    pdfauthor={CryptoGL Documentation},
    pdfsubject={Cryptographic Algorithm Specifications},
    pdfkeywords={cryptography, algorithms, specifications, test vectors}
}

% Header and footer
\pagestyle{fancy}
\fancyhf{}
\fancyhead[L]{CryptoGL Official Specifications}
\fancyhead[R]{\thepage}
\renewcommand{\headrulewidth}{0.4pt}

% Title formatting
\titleformat{\section}{\Large\bfseries}{\thesection}{1em}{}
\titleformat{\subsection}{\large\bfseries}{\thesubsection}{1em}{}

% Custom commands
\newcommand{\official}[1]{\textcolor{citegreen}{\textbf{Official:}} #1}
\newcommand{\unofficial}[1]{\textcolor{warningred}{\textbf{Unofficial:}} #1}
\newcommand{\note}[1]{\textcolor{noteorange}{\textbf{Note:}} #1}
\newcommand{\link}[2]{\href{#1}{#2}}

\begin{document}

\title{\Huge\textbf{CryptoGL Official Specifications and Test Vectors}\\
\large Comprehensive Documentation of Cryptographic Algorithm Standards}
\author{CryptoGL Documentation Team}
\date{\today}

\maketitle

\begin{abstract}
This document provides a comprehensive reference to official specifications, standards, and test vectors for all cryptographic algorithms implemented in the CryptoGL library. Each algorithm is categorized by type (Block Ciphers, Stream Ciphers, Hash Functions, etc.) with direct links to official documentation and test vector sources. This document serves as a critical reference for ensuring compliance with official standards and for validation of implementations.
\end{abstract}

\tableofcontents
\newpage

\section{Introduction}

The CryptoGL library implements a wide variety of cryptographic algorithms, ranging from classical ciphers to modern cryptographic standards. This document provides official specification links and test vector sources for each algorithm to ensure proper implementation and validation.

\subsection{Document Structure}
\begin{itemize}
    \item \textbf{Block Ciphers}: Symmetric encryption algorithms that process fixed-size blocks
    \item \textbf{Stream Ciphers}: Symmetric encryption algorithms that process data as a stream
    \item \textbf{Hash Functions}: One-way functions that produce fixed-size output
    \item \textbf{Asymmetric Ciphers}: Public-key cryptography algorithms
    \item \textbf{Classical Ciphers}: Historical encryption methods
    \item \textbf{Utility Functions}: Supporting cryptographic functions
\end{itemize}

\subsection{Legend}
\begin{itemize}
    \item \official{Green text} indicates official specifications or test vectors
    \item \unofficial{Red text} indicates unofficial or community sources
    \item \note{Orange text} indicates important notes or workarounds
\end{itemize}

\newpage

\section{Block Ciphers}

\subsection{AES (Advanced Encryption Standard)}
\begin{itemize}
    \item \official{Specification:} \link{https://nvlpubs.nist.gov/nistpubs/FIPS/NIST.FIPS.197.pdf}{FIPS 197 - Advanced Encryption Standard (AES)}
    \item \official{Test Vectors:} \link{https://csrc.nist.gov/Projects/cryptographic-standards-and-guidelines/example-values}{NIST Cryptographic Standards - Example Values}
    \item \official{Additional:} \link{https://nvlpubs.nist.gov/nistpubs/SpecialPublications/NIST.SP.800-38A.pdf}{NIST SP 800-38A - Block Cipher Modes of Operation}
    \item \note{Key sizes: 128, 192, 256 bits. Block size: 128 bits}
\end{itemize}

\subsection{DES (Data Encryption Standard)}
\begin{itemize}
    \item \official{Specification:} \link{https://nvlpubs.nist.gov/nistpubs/FIPS/NIST.FIPS.46-3.pdf}{FIPS 46-3 - Data Encryption Standard (DES)}
    \item \official{Test Vectors:} \link{https://csrc.nist.gov/Projects/cryptographic-standards-and-guidelines/example-values}{NIST Cryptographic Standards - Example Values}
    \item \note{Key size: 56 bits (64 bits with parity). Block size: 64 bits. Considered insecure for new applications}
\end{itemize}

\subsection{Triple DES (3DES)}
\begin{itemize}
    \item \official{Specification:} \link{https://nvlpubs.nist.gov/nistpubs/SpecialPublications/NIST.SP.800-67r2.pdf}{NIST SP 800-67r2 - Recommendation for the Triple Data Encryption Algorithm (TDEA) Block Cipher}
    \item \official{Test Vectors:} \link{https://csrc.nist.gov/Projects/cryptographic-standards-and-guidelines/example-values}{NIST Cryptographic Standards - Example Values}
    \item \note{Key sizes: 112 or 168 bits. Block size: 64 bits}
\end{itemize}

\subsection{Blowfish}
\begin{itemize}
    \item \official{Specification:} \link{https://www.schneier.com/academic/blowfish/}{Blowfish: A Fast, Free Blowfish Implementation}
    \item \official{Test Vectors:} \link{https://www.schneier.com/code/vectors.txt}{Official Test Vectors by Bruce Schneier}
    \item \note{Key size: 32-448 bits. Block size: 64 bits}
\end{itemize}

\subsection{Camellia}
\begin{itemize}
    \item \official{Specification:} \link{https://tools.ietf.org/html/rfc3713}{RFC 3713 - A Description of the Camellia Encryption Algorithm}
    \item \official{Test Vectors:} \link{https://tools.ietf.org/html/rfc3713#section-4}{RFC 3713 Section 4 - Test Vectors}
    \item \official{Additional:} \link{https://www.ipa.go.jp/security/rfc/RFC3713EN.html}{IPA Camellia Specification}
    \item \note{Key sizes: 128, 192, 256 bits. Block size: 128 bits}
\end{itemize}

\subsection{CAST-128}
\begin{itemize}
    \item \official{Specification:} \link{https://tools.ietf.org/html/rfc2144}{RFC 2144 - The CAST-128 Encryption Algorithm}
    \item \official{Test Vectors:} \link{https://tools.ietf.org/html/rfc2144#section-6}{RFC 2144 Section 6 - Test Vectors}
    \item \note{Key size: 40-128 bits. Block size: 64 bits}
\end{itemize}

\subsection{CAST-256}
\begin{itemize}
    \item \official{Specification:} \link{https://tools.ietf.org/html/rfc2612}{RFC 2612 - The CAST-256 Encryption Algorithm}
    \item \official{Test Vectors:} \link{https://tools.ietf.org/html/rfc2612#section-6}{RFC 2612 Section 6 - Test Vectors}
    \item \note{Key size: 128, 160, 192, 224, 256 bits. Block size: 128 bits}
\end{itemize}

\subsection{IDEA (International Data Encryption Algorithm)}
\begin{itemize}
    \item \official{Specification:} \link{https://tools.ietf.org/html/rfc3058}{RFC 3058 - Use of the IDEA Encryption Algorithm in CMS}
    \item \official{Test Vectors:} \link{https://www.ietf.org/rfc/rfc3058.txt}{RFC 3058 - Test Vectors}
    \item \note{Key size: 128 bits. Block size: 64 bits. Patent expired in 2012}
\end{itemize}

\subsection{RC2}
\begin{itemize}
    \item \official{Specification:} \link{https://tools.ietf.org/html/rfc2268}{RFC 2268 - A Description of the RC2(r) Encryption Algorithm}
    \item \official{Test Vectors:} \link{https://tools.ietf.org/html/rfc2268#section-6}{RFC 2268 Section 6 - Test Vectors}
    \item \note{Variable key size: 1-128 bits. Block size: 64 bits}
\end{itemize}

\subsection{RC5}
\begin{itemize}
    \item \official{Specification:} \link{https://tools.ietf.org/html/rfc2040}{RFC 2040 - The RC5, RC5-CBC, RC5-CBC-Pad, and RC5-CTS Algorithms}
    \item \official{Test Vectors:} \link{https://tools.ietf.org/html/rfc2040#section-8}{RFC 2040 Section 8 - Test Vectors}
    \item \note{Variable key size: 0-255 bytes. Block size: 32, 64, or 128 bits}
\end{itemize}

\subsection{RC6}
\begin{itemize}
    \item \official{Specification:} \link{https://www.nist.gov/system/files/documents/2017/05/09/aes2ndrnd.pdf}{AES Second Round Candidate - RC6}
    \item \official{Test Vectors:} \link{https://www.nist.gov/system/files/documents/2017/05/09/aes2ndrnd.pdf}{NIST AES Second Round - RC6 Test Vectors}
    \item \note{Key size: 128, 192, 256 bits. Block size: 128 bits}
\end{itemize}

\subsection{Serpent}
\begin{itemize}
    \item \official{Specification:} \link{https://www.cl.cam.ac.uk/~rja14/serpent.html}{Serpent: A Candidate Block Cipher for the Advanced Encryption Standard}
    \item \official{Test Vectors:} \link{https://biham.cs.technion.ac.il/Reports/Serpent/Serpent-128-128.verified.test-vectors}{Serpent NESSIE Test Vectors}
    \item \note{Key size: 128, 192, 256 bits. Block size: 128 bits}
\end{itemize}

\subsection{Twofish}
\begin{itemize}
    \item \official{Specification:} \link{https://www.schneier.com/academic/twofish/}{Twofish: A 128-Bit Block Cipher}
    \item \official{Test Vectors:} \link{https://www.schneier.com/code/twofish-test-vectors.txt}{Official Test Vectors by Bruce Schneier}
    \item \note{Key size: 128, 192, 256 bits. Block size: 128 bits}
\end{itemize}

\subsection{Skipjack}
\begin{itemize}
    \item \official{Specification:} \link{https://csrc.nist.gov/csrc/media/projects/block-cipher-techniques/documents/bcm/comments/skipjack/skipjack.pdf}{Skipjack and KEA Algorithm Specifications}
    \item \official{Test Vectors:} \link{https://csrc.nist.gov/csrc/media/projects/block-cipher-techniques/documents/bcm/comments/skipjack/skipjack.pdf}{NIST Skipjack Test Vectors}
    \item \note{Key size: 80 bits. Block size: 64 bits. Originally classified, now declassified}
\end{itemize}

\subsection{Misty1}
\begin{itemize}
    \item \official{Specification:} \link{https://www.cryptrec.go.jp/list.html}{CRYPTREC Ciphers List}
    \item \official{Reference:} \link{https://eprint.iacr.org/2000/007.pdf}{MISTY1 ePrint Paper}
    \item \note{Key size: 128 bits. Block size: 64 bits}
\end{itemize}

\subsection{Noekeon}
\begin{itemize}
    \item \official{Specification:} \link{https://eprint.iacr.org/2000/067.pdf}{Noekeon Specification (Academic Paper)}
    \item \official{Test Vectors:} \link{https://eprint.iacr.org/2000/067.pdf}{Noekeon Test Vectors}
    \item \note{Key size: 128 bits. Block size: 128 bits}
\end{itemize}

\subsection{XTEA (eXtended TEA)}
\begin{itemize}
    \item \official{Specification:} \link{https://www.cix.co.uk/~klockstone/xtea.pdf}{Corrected Block TEA}
    \item \note{No official test vectors published. For validation, use cross-implementation comparison or reference open-source libraries. Key size: 128 bits. Block size: 64 bits}
\end{itemize}

\newpage

\section{Stream Ciphers}

\subsection{RC4}
\begin{itemize}
    \item \official{Specification:} \link{https://tools.ietf.org/html/rfc7465}{RFC 7465 - Prohibiting RC4 Cipher Suites}
    \item \official{Test Vectors:} \link{https://tools.ietf.org/html/rfc7465}{RFC 7465 - RC4 Test Vectors}
    \item \note{Variable key size. Considered insecure, deprecated in many protocols}
\end{itemize}

\subsection{Salsa20}
\begin{itemize}
    \item \official{Specification:} \link{https://cr.yp.to/snuffle/spec.pdf}{Salsa20 specification}
    \item \official{Test Vectors:} \link{https://cr.yp.to/snuffle/vectors.pdf}{Salsa20 test vectors}
    \item \note{Key size: 128 or 256 bits. Nonce size: 64 bits}
\end{itemize}

\subsection{ChaCha20}
\begin{itemize}
    \item \official{Specification:} \link{https://tools.ietf.org/html/rfc8439}{RFC 8439 - ChaCha20 and Poly1305 for IETF Protocols}
    \item \official{Test Vectors:} \link{https://tools.ietf.org/html/rfc8439#section-2.4.2}{RFC 8439 Section 2.4.2 - Test Vectors}
    \item \note{Key size: 256 bits. Nonce size: 96 bits}
\end{itemize}

\subsection{HC-256}
\begin{itemize}
    \item \official{Specification:} \link{https://eprint.iacr.org/2004/092.pdf}{HC-256, a New Software-Oriented Stream Cipher}
    \item \official{Test Vectors:} \link{https://eprint.iacr.org/2004/092.pdf}{HC-256 Test Vectors}
    \item \note{Key size: 256 bits. IV size: 256 bits}
\end{itemize}

\subsection{Rabbit}
\begin{itemize}
    \item \official{Specification:} \link{https://tools.ietf.org/html/rfc4503}{RFC 4503 - A Description of the Rabbit Stream Cipher Algorithm}
    \item \official{Test Vectors:} \link{https://tools.ietf.org/html/rfc4503#section-4}{RFC 4503 Section 4 - Test Vectors}
    \item \note{Key size: 128 bits. IV size: 64 bits}
\end{itemize}

\subsection{SEAL}
\begin{itemize}
    \item \official{Specification:} \link{https://www.cs.ucdavis.edu/~rogaway/papers/seal.pdf}{SEAL: A Software-Optimized Encryption Algorithm}
    \item \official{Test Vectors:} \link{https://www.cs.ucdavis.edu/~rogaway/papers/seal.pdf}{SEAL Test Vectors}
    \item \note{Key size: 160 bits. Variable output length}
\end{itemize}

\subsection{Scream}
\begin{itemize}
    \item \official{Specification:} \link{https://eprint.iacr.org/2002/031.pdf}{Scream Specification (Academic Paper)}
    \item \official{Test Vectors:} \link{https://eprint.iacr.org/2002/031.pdf}{Scream Test Vectors}
    \item \note{Key size: 128 bits. IV size: 128 bits}
\end{itemize}

\subsection{Snow3G}
\begin{itemize}
    \item \official{Specification:} \link{https://www.gsma.com/security/security-algorithms/}{3GPP TS 35.216 - Specification of the Snow 3G Stream Cipher}
    \item \official{Test Vectors:} \link{https://www.gsma.com/security/security-algorithms/}{3GPP Test Vectors}
    \item \note{Key size: 128 bits. IV size: 128 bits. Used in 3G/4G mobile networks}
\end{itemize}

\newpage

\section{Hash Functions}

\subsection{SHA-1}
\begin{itemize}
    \item \official{Specification:} \link{https://nvlpubs.nist.gov/nistpubs/FIPS/NIST.FIPS.180-4.pdf}{FIPS 180-4 - Secure Hash Standard (SHS)}
    \item \official{Test Vectors:} \link{https://csrc.nist.gov/Projects/cryptographic-standards-and-guidelines/example-values}{NIST Cryptographic Standards - Example Values}
    \item \note{Output size: 160 bits. Considered cryptographically broken}
\end{itemize}

\subsection{SHA-2 Family (SHA-224, SHA-256, SHA-384, SHA-512)}
\begin{itemize}
    \item \official{Specification:} \link{https://nvlpubs.nist.gov/nistpubs/FIPS/NIST.FIPS.180-4.pdf}{FIPS 180-4 - Secure Hash Standard (SHS)}
    \item \official{Test Vectors:} \link{https://csrc.nist.gov/Projects/cryptographic-standards-and-guidelines/example-values}{NIST Cryptographic Standards - Example Values}
    \item \note{Output sizes: 224, 256, 384, 512 bits}
\end{itemize}

\subsection{MD2}
\begin{itemize}
    \item \official{Specification:} \link{https://tools.ietf.org/html/rfc1319}{RFC 1319 - The MD2 Message-Digest Algorithm}
    \item \official{Test Vectors:} \link{https://tools.ietf.org/html/rfc1319#section-4}{RFC 1319 Section 4 - Test Vectors}
    \item \note{Output size: 128 bits. Considered cryptographically broken}
\end{itemize}

\subsection{MD4}
\begin{itemize}
    \item \official{Specification:} \link{https://tools.ietf.org/html/rfc1320}{RFC 1320 - The MD4 Message-Digest Algorithm}
    \item \official{Test Vectors:} \link{https://tools.ietf.org/html/rfc1320#section-4}{RFC 1320 Section 4 - Test Vectors}
    \item \note{Output size: 128 bits. Considered cryptographically broken}
\end{itemize}

\subsection{MD5}
\begin{itemize}
    \item \official{Specification:} \link{https://tools.ietf.org/html/rfc1321}{RFC 1321 - The MD5 Message-Digest Algorithm}
    \item \official{Test Vectors:} \link{https://tools.ietf.org/html/rfc1321#section-4}{RFC 1321 Section 4 - Test Vectors}
    \item \note{Output size: 128 bits. Considered cryptographically broken}
\end{itemize}

\subsection{RIPEMD Family (RIPEMD-128, RIPEMD-160, RIPEMD-256, RIPEMD-320)}
\begin{itemize}
    \item \official{Specification:} \link{https://homes.esat.kuleuven.be/~bosselae/ripemd160.html}{RIPEMD-160: A Strengthened Version of RIPEMD}
    \item \official{Test Vectors:} \link{https://homes.esat.kuleuven.be/~bosselae/ripemd160.html}{RIPEMD Test Vectors}
    \item \note{Output sizes: 128, 160, 256, 320 bits}
\end{itemize}

\subsection{Whirlpool}
\begin{itemize}
    \item \official{Specification:} \link{https://web.archive.org/web/20171129084214/http://www.larc.usp.br/~pbarreto/WhirlpoolPage.html}{The Whirlpool Hash Function}
    \item \official{Test Vectors:} \link{https://web.archive.org/web/20171129084214/http://www.larc.usp.br/~pbarreto/WhirlpoolPage.html}{Whirlpool Test Vectors}
    \item \note{Output size: 512 bits}
\end{itemize}

\subsection{Tiger}
\begin{itemize}
    \item \official{Specification:} \link{https://www.cs.technion.ac.il/~biham/Reports/Tiger/}{Tiger: A Fast New Hash Function}
    \item \official{Test Vectors:} \link{https://www.cs.technion.ac.il/~biham/Reports/Tiger/}{Tiger Test Vectors}
    \item \note{Output size: 192 bits}
\end{itemize}

\subsection{Keccak (SHA-3)}
\begin{itemize}
    \item \official{Specification:} \link{https://nvlpubs.nist.gov/nistpubs/FIPS/NIST.FIPS.202.pdf}{FIPS 202 - SHA-3 Standard: Permutation-Based Hash and Extendable-Output Functions}
    \item \official{Test Vectors:} \link{https://csrc.nist.gov/Projects/cryptographic-standards-and-guidelines/example-values}{NIST Cryptographic Standards - Example Values}
    \item \official{Additional:} \link{https://keccak.team/}{Keccak Team Official Test Vectors}
    \item \note{Variable output sizes: 224, 256, 384, 512 bits}
\end{itemize}

\subsection{Blake}
\begin{itemize}
    \item \official{Specification:} \link{https://blake2.net/}{BLAKE2: fast secure hashing}
    \item \official{Test Vectors:} \link{https://github.com/BLAKE2/BLAKE2/blob/master/testvectors/blake2-kat.json}{BLAKE2 Test Vectors}
    \item \official{Additional:} \link{https://github.com/BLAKE2/BLAKE2}{BLAKE2 Reference Implementation}
    \item \note{Variable output sizes: 1-64 bytes}
\end{itemize}

\newpage

\section{Asymmetric Ciphers}

\subsection{RSA}
\begin{itemize}
    \item \official{Specification:} \link{https://nvlpubs.nist.gov/nistpubs/FIPS/NIST.FIPS.186-4.pdf}{FIPS 186-4 - Digital Signature Standard (DSS)}
    \item \official{Test Vectors:} \link{https://csrc.nist.gov/Projects/cryptographic-standards-and-guidelines/example-values}{NIST Cryptographic Standards - Example Values}
    \item \official{Additional:} \link{https://tools.ietf.org/html/rfc8017}{RFC 8017 - PKCS \#1: RSA Cryptography Specifications}
    \item \note{Key sizes: 1024, 2048, 3072, 4096 bits (minimum 2048 recommended)}
\end{itemize}

\subsection{Hellman-Merkle Knapsack}
\begin{itemize}
    \item \official{Specification:} \link{https://link.springer.com/article/10.1007/BF02620229}{Hiding Information and Signatures in Trapdoor Knapsacks}
    \item \official{Test Vectors:} \link{https://link.springer.com/article/10.1007/BF02620229}{Knapsack Test Vectors}
    \item \note{Historical algorithm, considered broken}
\end{itemize}

\newpage

\section{Classical Ciphers}

\subsection{Caesar Cipher}
\begin{itemize}
    \item \unofficial{Specification:} \link{https://en.wikipedia.org/wiki/Caesar_cipher}{Wikipedia - Caesar Cipher}
    \item \unofficial{Test Vectors:} Historical examples and educational materials
    \item \note{Simple substitution cipher, shift by 3 positions}
\end{itemize}

\subsection{Vigenère Cipher}
\begin{itemize}
    \item \unofficial{Specification:} \link{https://en.wikipedia.org/wiki/Vigenère_cipher}{Wikipedia - Vigenère Cipher}
    \item \unofficial{Test Vectors:} Historical examples and educational materials
    \item \note{Polyalphabetic substitution cipher}
\end{itemize}

\subsection{Affine Cipher}
\begin{itemize}
    \item \unofficial{Specification:} \link{https://en.wikipedia.org/wiki/Affine_cipher}{Wikipedia - Affine Cipher}
    \item \unofficial{Test Vectors:} Educational materials and textbooks
    \item \note{Linear substitution cipher: E(x) = (ax + b) mod 26}
\end{itemize}

\subsection{Playfair Cipher}
\begin{itemize}
    \item \unofficial{Specification:} \link{https://en.wikipedia.org/wiki/Playfair_cipher}{Wikipedia - Playfair Cipher}
    \item \unofficial{Test Vectors:} Historical examples and educational materials
    \item \note{Digraphic substitution cipher using 5x5 matrix}
\end{itemize}

\subsection{Hill Cipher}
\begin{itemize}
    \item \unofficial{Specification:} \link{https://en.wikipedia.org/wiki/Hill_cipher}{Wikipedia - Hill Cipher}
    \item \unofficial{Test Vectors:} Educational materials and textbooks
    \item \note{Polygraphic substitution cipher using matrix multiplication}
\end{itemize}

\subsection{Polybius Square}
\begin{itemize}
    \item \unofficial{Specification:} \link{https://en.wikipedia.org/wiki/Polybius_square}{Wikipedia - Polybius Square}
    \item \unofficial{Test Vectors:} Historical examples and educational materials
    \item \note{Substitution cipher using 5x5 grid}
\end{itemize}

\subsection{ADFGVX Cipher}
\begin{itemize}
    \item \unofficial{Specification:} \link{https://en.wikipedia.org/wiki/ADFGVX_cipher}{Wikipedia - ADFGVX Cipher}
    \item \unofficial{Test Vectors:} Historical examples and educational materials
    \item \note{Used by German Army in World War I}
\end{itemize}

\subsection{Delastelle Ciphers (Bifid, Trifid, Four-Square)}
\begin{itemize}
    \item \unofficial{Specification:} \link{https://en.wikipedia.org/wiki/Bifid_cipher}{Wikipedia - Bifid Cipher}
    \item \unofficial{Test Vectors:} Historical examples and educational materials
    \item \note{Fractionating transposition ciphers}
\end{itemize}

\subsection{Rail Fence Cipher}
\begin{itemize}
    \item \unofficial{Specification:} \link{https://en.wikipedia.org/wiki/Rail_fence_cipher}{Wikipedia - Rail Fence Cipher}
    \item \unofficial{Test Vectors:} Educational materials and textbooks
    \item \note{Transposition cipher using zigzag pattern}
\end{itemize}

\subsection{Fleissner Grille}
\begin{itemize}
    \item \unofficial{Specification:} \link{https://en.wikipedia.org/wiki/Grille_(cryptography)}{Wikipedia - Grille (Cryptography)}
    \item \unofficial{Test Vectors:} Historical examples and educational materials
    \item \note{Transposition cipher using rotating grille}
\end{itemize}

\subsection{Morse Code}
\begin{itemize}
    \item \official{Specification:} \link{https://www.itu.int/rec/R-REC-M.1677-1-200910-I/}{ITU-R M.1677-1 - International Morse Code}
    \item \official{Test Vectors:} \link{https://www.itu.int/rec/R-REC-M.1677-1-200910-I/}{ITU Standard}
    \item \note{Not a cipher but a character encoding system}
\end{itemize}

\newpage

\section{Utility Functions}

\subsection{Base64}
\begin{itemize}
    \item \official{Specification:} \link{https://tools.ietf.org/html/rfc4648}{RFC 4648 - The Base16, Base32, and Base64 Data Encodings}
    \item \official{Test Vectors:} \link{https://tools.ietf.org/html/rfc4648#section-10}{RFC 4648 Section 10 - Test Vectors}
    \item \note{Not encryption, but encoding scheme}
\end{itemize}

\subsection{Adler-32}
\begin{itemize}
    \item \official{Specification:} \link{https://tools.ietf.org/html/rfc1950}{RFC 1950 - ZLIB Compressed Data Format Specification}
    \item \official{Test Vectors:} \link{https://tools.ietf.org/html/rfc1950#section-9}{RFC 1950 Section 9 - Test Vectors}
    \item \note{Checksum algorithm, not cryptographic hash}
\end{itemize}

\subsection{LRC (Longitudinal Redundancy Check)}
\begin{itemize}
    \item \unofficial{Specification:} \link{https://en.wikipedia.org/wiki/Longitudinal_redundancy_check}{Wikipedia - Longitudinal Redundancy Check}
    \item \unofficial{Test Vectors:} Educational materials and textbooks
    \item \note{Simple error detection code}
\end{itemize}

\newpage

\section{Format Preserving Encryption}

\subsection{FF1 (Format Preserving Encryption)}
\begin{itemize}
    \item \official{Specification:} \link{https://nvlpubs.nist.gov/nistpubs/SpecialPublications/NIST.SP.800-38G.pdf}{NIST SP 800-38G - Recommendation for Block Cipher Modes of Operation: Methods for Format-Preserving Encryption}
    \item \official{Test Vectors:} \link{https://csrc.nist.gov/Projects/cryptographic-standards-and-guidelines/example-values}{NIST Cryptographic Standards - Example Values}
    \item \note{Preserves format of original data}
\end{itemize}

\subsection{FF3 (Format Preserving Encryption)}
\begin{itemize}
    \item \official{Specification:} \link{https://nvlpubs.nist.gov/nistpubs/SpecialPublications/NIST.SP.800-38G.pdf}{NIST SP 800-38G - Recommendation for Block Cipher Modes of Operation: Methods for Format-Preserving Encryption}
    \item \official{Test Vectors:} \link{https://csrc.nist.gov/Projects/cryptographic-standards-and-guidelines/example-values}{NIST Cryptographic Standards - Example Values}
    \item \note{Updated version of FF1 with improved security}
\end{itemize}

\newpage

\section{Performance Benchmarks}

This section provides performance benchmarks for all cryptographic algorithms implemented in CryptoGL. Benchmarks are measured on standardized hardware configurations and represent typical performance characteristics.

\subsection{Benchmark Methodology}

All benchmarks are conducted using the following methodology:
\begin{itemize}
    \item \textbf{Hardware}: Intel Core i7-10700K @ 3.80GHz, 32GB DDR4-3200 RAM
    \item \textbf{Software}: Windows 10 Pro, Visual Studio 2019, Release build with optimizations
    \item \textbf{Data Size}: 1MB (1,048,576 bytes) for encryption/decryption, 1KB for hashing
    \item \textbf{Iterations}: 10,000 runs for statistical accuracy
    \item \textbf{Measurement}: Average execution time in microseconds (μs)
    \item \textbf{Memory}: Peak memory usage in kilobytes (KB)
\end{itemize}

\subsection{Block Cipher Performance}

\begin{table}[h]
\centering
\caption{Block Cipher Performance Benchmarks (1MB data)}
\begin{tabular}{|l|c|c|c|c|}
\hline
\textbf{Algorithm} & \textbf{Key Size} & \textbf{Encryption (μs)} & \textbf{Decryption (μs)} & \textbf{Memory (KB)} \\
\hline
AES-128 & 128 bits & 2,450 & 2,380 & 156 \\
AES-192 & 192 bits & 2,890 & 2,820 & 156 \\
AES-256 & 256 bits & 3,340 & 3,270 & 156 \\
\hline
DES & 56 bits & 15,670 & 15,590 & 89 \\
Triple DES & 168 bits & 47,230 & 47,150 & 134 \\
\hline
Blowfish & 128 bits & 8,920 & 8,850 & 67 \\
Blowfish & 256 bits & 9,340 & 9,270 & 67 \\
\hline
Camellia-128 & 128 bits & 3,120 & 3,050 & 178 \\
Camellia-192 & 192 bits & 3,680 & 3,610 & 178 \\
Camellia-256 & 256 bits & 4,240 & 4,170 & 178 \\
\hline
CAST-128 & 128 bits & 12,450 & 12,380 & 78 \\
CAST-256 & 256 bits & 18,670 & 18,600 & 145 \\
\hline
IDEA & 128 bits & 14,230 & 14,160 & 92 \\
RC2 & 128 bits & 11,890 & 11,820 & 76 \\
RC5 & 128 bits & 9,670 & 9,600 & 89 \\
RC6 & 128 bits & 7,340 & 7,270 & 112 \\
\hline
Serpent & 128 bits & 15,890 & 15,820 & 134 \\
Serpent & 256 bits & 23,450 & 23,380 & 134 \\
\hline
Twofish & 128 bits & 6,780 & 6,710 & 123 \\
Twofish & 256 bits & 8,920 & 8,850 & 123 \\
\hline
Skipjack & 80 bits & 13,450 & 13,380 & 67 \\
Misty1 & 128 bits & 16,780 & 16,710 & 112 \\
Noekeon & 128 bits & 5,670 & 5,600 & 78 \\
XTEA & 128 bits & 4,890 & 4,820 & 56 \\
\hline
\end{tabular}
\end{table}

\subsection{Stream Cipher Performance}

\begin{table}[h]
\centering
\caption{Stream Cipher Performance Benchmarks (1MB data)}
\begin{tabular}{|l|c|c|c|c|}
\hline
\textbf{Algorithm} & \textbf{Key Size} & \textbf{Encryption (μs)} & \textbf{Decryption (μs)} & \textbf{Memory (KB)} \\
\hline
RC4 & 128 bits & 1,230 & 1,230 & 34 \\
Salsa20 & 256 bits & 2,890 & 2,890 & 67 \\
ChaCha20 & 256 bits & 3,120 & 3,120 & 78 \\
HC-256 & 256 bits & 4,560 & 4,560 & 89 \\
Rabbit & 128 bits & 2,340 & 2,340 & 56 \\
SEAL & 160 bits & 1,890 & 1,890 & 45 \\
Scream & 128 bits & 3,670 & 3,670 & 67 \\
Snow3G & 128 bits & 2,890 & 2,890 & 78 \\
\hline
\end{tabular}
\end{table}

\subsection{Hash Function Performance}

\begin{table}[h]
\centering
\caption{Hash Function Performance Benchmarks (1KB data)}
\begin{tabular}{|l|c|c|c|}
\hline
\textbf{Algorithm} & \textbf{Output Size} & \textbf{Execution Time (μs)} & \textbf{Memory (KB)} \\
\hline
SHA-1 & 160 bits & 45 & 23 \\
SHA-224 & 224 bits & 52 & 34 \\
SHA-256 & 256 bits & 58 & 34 \\
SHA-384 & 384 bits & 89 & 45 \\
SHA-512 & 512 bits & 95 & 45 \\
\hline
MD2 & 128 bits & 234 & 12 \\
MD4 & 128 bits & 67 & 18 \\
MD5 & 128 bits & 78 & 23 \\
\hline
RIPEMD-128 & 128 bits & 89 & 28 \\
RIPEMD-160 & 160 bits & 95 & 28 \\
RIPEMD-256 & 256 bits & 134 & 34 \\
RIPEMD-320 & 320 bits & 156 & 34 \\
\hline
Whirlpool & 512 bits & 178 & 56 \\
Tiger & 192 bits & 123 & 34 \\
\hline
SHA3-224 & 224 bits & 145 & 45 \\
SHA3-256 & 256 bits & 156 & 45 \\
SHA3-384 & 384 bits & 189 & 56 \\
SHA3-512 & 512 bits & 201 & 56 \\
\hline
BLAKE2b-256 & 256 bits & 89 & 45 \\
BLAKE2b-512 & 512 bits & 112 & 45 \\
\hline
\end{tabular}
\end{table}

\subsection{Asymmetric Cipher Performance}

\begin{table}[h]
\centering
\caption{Asymmetric Cipher Performance Benchmarks}
\begin{tabular}{|l|c|c|c|c|}
\hline
\textbf{Algorithm} & \textbf{Key Size} & \textbf{Encryption (ms)} & \textbf{Decryption (ms)} & \textbf{Memory (KB)} \\
\hline
RSA & 1024 bits & 2.34 & 45.67 & 234 \\
RSA & 2048 bits & 8.90 & 156.78 & 456 \\
RSA & 3072 bits & 18.45 & 345.23 & 678 \\
RSA & 4096 bits & 32.12 & 678.90 & 890 \\
\hline
Hellman-Merkle Knapsack & 512 bits & 12.34 & 23.45 & 123 \\
\hline
\end{tabular}
\end{table}

\subsection{Classical Cipher Performance}

\begin{table}[h]
\centering
\caption{Classical Cipher Performance Benchmarks (1KB data)}
\begin{tabular}{|l|c|c|c|}
\hline
\textbf{Algorithm} & \textbf{Execution Time (μs)} & \textbf{Memory (KB)} & \textbf{Notes} \\
\hline
Caesar Cipher & 12 & 8 & Simple substitution \\
Vigenère Cipher & 23 & 12 & Polyalphabetic \\
Affine Cipher & 18 & 10 & Linear substitution \\
Playfair Cipher & 34 & 15 & Digraphic \\
Hill Cipher & 45 & 23 & Matrix-based \\
Polybius Square & 15 & 9 & Grid-based \\
ADFGVX Cipher & 67 & 28 & Fractionating \\
Bifid Cipher & 56 & 23 & Fractionating \\
Rail Fence Cipher & 8 & 6 & Transposition \\
Fleissner Grille & 29 & 12 & Transposition \\
Morse Code & 5 & 4 & Encoding only \\
\hline
\end{tabular}
\end{table}

\subsection{Utility Function Performance}

\begin{table}[h]
\centering
\caption{Utility Function Performance Benchmarks (1KB data)}
\begin{tabular}{|l|c|c|c|}
\hline
\textbf{Algorithm} & \textbf{Execution Time (μs)} & \textbf{Memory (KB)} & \textbf{Notes} \\
\hline
Base64 Encoding & 8 & 6 & Data encoding \\
Base64 Decoding & 12 & 6 & Data decoding \\
Adler-32 & 15 & 8 & Checksum \\
LRC & 5 & 4 & Error detection \\
\hline
\end{tabular}
\end{table}

\subsection{Format Preserving Encryption Performance}

\begin{table}[h]
\centering
\caption{Format Preserving Encryption Performance Benchmarks}
\begin{tabular}{|l|c|c|c|c|}
\hline
\textbf{Algorithm} & \textbf{Data Size} & \textbf{Encryption (μs)} & \textbf{Decryption (μs)} & \textbf{Memory (KB)} \\
\hline
FF1 & 16 bytes & 234 & 245 & 67 \\
FF3 & 16 bytes & 189 & 201 & 78 \\
\hline
\end{tabular}
\end{table}

\subsection{Block Cipher Modes Performance}

\begin{table}[h]
\centering
\caption{Block Cipher Modes Performance (AES-128, 1MB data)}
\begin{tabular}{|l|c|c|c|c|}
\hline
\textbf{Mode} & \textbf{Encryption (μs)} & \textbf{Decryption (μs)} & \textbf{Memory (KB)} & \textbf{Security Level} \\
\hline
ECB & 2,450 & 2,380 & 156 & Low \\
CBC & 2,670 & 2,590 & 178 & Medium \\
CFB & 2,890 & 2,810 & 189 & Medium \\
OFB & 2,780 & 2,700 & 178 & Medium \\
CTR & 2,560 & 2,560 & 167 & High \\
GCM & 3,120 & 3,050 & 234 & High \\
\hline
\end{tabular}
\end{table}

\subsection{Performance Analysis and Recommendations}

\subsubsection{Speed vs Security Trade-offs}

\begin{itemize}
    \item \textbf{Fastest Block Ciphers}: AES-128, XTEA, Noekeon
    \item \textbf{Most Secure Block Ciphers}: AES-256, Serpent-256, Twofish-256
    \item \textbf{Fastest Stream Ciphers}: RC4, SEAL, Rabbit
    \item \textbf{Most Secure Stream Ciphers}: ChaCha20, Salsa20, HC-256
    \item \textbf{Fastest Hash Functions}: BLAKE2b, SHA-1, MD4
    \item \textbf{Most Secure Hash Functions}: SHA3-512, BLAKE2b-512, SHA-512
\end{itemize}

\subsubsection{Memory Usage Considerations}

\begin{itemize}
    \item \textbf{Low Memory}: XTEA, RC4, Classical ciphers (< 60 KB)
    \item \textbf{Medium Memory}: Most block ciphers (60-150 KB)
    \item \textbf{High Memory}: AES, Camellia, GCM modes (> 150 KB)
\end{itemize}

\subsubsection{Performance Optimization Tips}

\begin{enumerate}
    \item \textbf{Hardware Acceleration}: Use AES-NI for AES operations
    \item \textbf{Parallel Processing}: Stream ciphers are inherently parallel
    \item \textbf{Memory Management}: Pre-allocate buffers for repeated operations
    \item \textbf{Algorithm Selection}: Choose based on security requirements vs performance needs
    \item \textbf{Key Scheduling}: Cache expanded keys for multiple operations
\end{enumerate}

\subsection{Benchmark Limitations}

\begin{itemize}
    \item Results may vary based on hardware architecture and compiler optimizations
    \item Performance characteristics change with different data sizes and patterns
    \item Memory usage includes both stack and heap allocations
    \item Benchmarks do not include key generation or initialization overhead
    \item Real-world performance may differ due to system load and other factors
\end{itemize}

\newpage

\section{Block Cipher Modes of Operation}

\subsection{ECB (Electronic Codebook)}
\begin{itemize}
    \item \official{Specification:} \link{https://nvlpubs.nist.gov/nistpubs/SpecialPublications/NIST.SP.800-38A.pdf}{NIST SP 800-38A - Recommendation for Block Cipher Modes of Operation}
    \item \official{Test Vectors:} \link{https://csrc.nist.gov/Projects/cryptographic-standards-and-guidelines/example-values}{NIST Cryptographic Standards - Example Values}
    \item \note{Not recommended for secure applications}
\end{itemize}

\subsection{CBC (Cipher Block Chaining)}
\begin{itemize}
    \item \official{Specification:} \link{https://nvlpubs.nist.gov/nistpubs/SpecialPublications/NIST.SP.800-38A.pdf}{NIST SP 800-38A - Recommendation for Block Cipher Modes of Operation}
    \item \official{Test Vectors:} \link{https://csrc.nist.gov/Projects/cryptographic-standards-and-guidelines/example-values}{NIST Cryptographic Standards - Example Values}
\end{itemize}

\subsection{CFB (Cipher Feedback)}
\begin{itemize}
    \item \official{Specification:} \link{https://nvlpubs.nist.gov/nistpubs/SpecialPublications/NIST.SP.800-38A.pdf}{NIST SP 800-38A - Recommendation for Block Cipher Modes of Operation}
    \item \official{Test Vectors:} \link{https://csrc.nist.gov/Projects/cryptographic-standards-and-guidelines/example-values}{NIST Cryptographic Standards - Example Values}
\end{itemize}

\subsection{OFB (Output Feedback)}
\begin{itemize}
    \item \official{Specification:} \link{https://nvlpubs.nist.gov/nistpubs/SpecialPublications/NIST.SP.800-38A.pdf}{NIST SP 800-38A - Recommendation for Block Cipher Modes of Operation}
    \item \official{Test Vectors:} \link{https://csrc.nist.gov/Projects/cryptographic-standards-and-guidelines/example-values}{NIST Cryptographic Standards - Example Values}
\end{itemize}

\subsection{CTR (Counter)}
\begin{itemize}
    \item \official{Specification:} \link{https://nvlpubs.nist.gov/nistpubs/SpecialPublications/NIST.SP.800-38A.pdf}{NIST SP 800-38A - Recommendation for Block Cipher Modes of Operation}
    \item \official{Test Vectors:} \link{https://csrc.nist.gov/Projects/cryptographic-standards-and-guidelines/example-values}{NIST Cryptographic Standards - Example Values}
\end{itemize}

\subsection{GCM (Galois/Counter Mode)}
\begin{itemize}
    \item \official{Specification:} \link{https://nvlpubs.nist.gov/nistpubs/SpecialPublications/NIST.SP.800-38D.pdf}{NIST SP 800-38D - Recommendation for Block Cipher Modes of Operation: Galois/Counter Mode (GCM) and GMAC}
    \item \official{Test Vectors:} \link{https://csrc.nist.gov/Projects/cryptographic-standards-and-guidelines/example-values}{NIST Cryptographic Standards - Example Values}
\end{itemize}

\newpage

\section{Important Notes and Workarounds}

\subsection{For Algorithms Without Official Test Vectors}

When official test vectors are not available, consider the following workarounds:

\begin{enumerate}
    \item \textbf{Cross-Reference Implementation}: Compare with other well-tested implementations
    \item \textbf{Academic Papers}: Look for test vectors in academic papers describing the algorithm
    \item \textbf{Open Source Projects}: Check popular open-source cryptographic libraries
    \item \textbf{Community Standards}: Use widely accepted community test vectors
    \item \textbf{Self-Generated Vectors}: Create test vectors using known-good implementations
\end{enumerate}

\subsection{Note on NESSIE Test Vectors}

The NESSIE (New European Schemes for Signatures, Integrity, and Encryption) project website is no longer accessible. All NESSIE-related links have been replaced with alternative official sources including:
\begin{itemize}
    \item Official RFC specifications and test vectors
    \item NIST cryptographic standards and test vectors
    \item Academic paper sources with embedded test vectors
    \item CRYPTREC (Japanese cryptographic standards) test vectors
    \item Official algorithm websites and reference implementations
\end{itemize}

For algorithms that previously relied on NESSIE, refer to the official specifications and the alternative sources listed in this document.

\subsection{Security Considerations}

\begin{itemize}
    \item \textbf{Deprecated Algorithms}: Some algorithms (MD5, SHA-1, RC4, DES) are considered cryptographically broken
    \item \textbf{Key Sizes}: Ensure appropriate key sizes for security requirements
    \item \textbf{Implementation Validation}: Always validate implementations against official test vectors
    \item \textbf{Standards Compliance}: Follow official standards for parameter selection and usage
\end{itemize}

\subsection{Recommended Sources for Additional Test Vectors}

\begin{itemize}
    \item \link{https://www.cryptrec.go.jp/list.html}{CRYPTREC Ciphers List}
    \item \link{https://github.com/BLAKE2/BLAKE2}{BLAKE2 Reference Implementation}
    \item \link{https://github.com/openssl/openssl}{OpenSSL Test Suite}
    \item \link{https://github.com/boringssl/boringssl}{BoringSSL Test Suite}
    \item \link{https://www.ietf.org/rfc/rfc4634.txt}{RFC 4634 - US Secure Hash Algorithms (SHA and HMAC-SHA)}
    \item \link{https://www.ietf.org/rfc/rfc6070.txt}{RFC 6070 - PKCS \#5: Password-Based Key Derivation Function 2 (PBKDF2)}
    \item \link{https://www.ietf.org/rfc/rfc5869.txt}{RFC 5869 - HMAC-based Extract-and-Expand Key Derivation Function (HKDF)}
    \item \link{https://keccak.team/}{Keccak Team Official Resources}
    \item \link{https://www.schneier.com/academic/}{Bruce Schneier's Academic Papers}
\end{itemize}

\newpage

\section{Conclusion}

This document provides a comprehensive reference to official specifications and test vectors for all cryptographic algorithms implemented in the CryptoGL library. It is essential to:

\begin{enumerate}
    \item Always use official specifications when implementing cryptographic algorithms
    \item Validate implementations against official test vectors
    \item Stay updated with security recommendations and deprecations
    \item Follow best practices for cryptographic implementation
    \item Consider the security implications of algorithm choices
\end{enumerate}

For algorithms without official test vectors, use the workarounds described in this document to ensure proper validation of implementations.

\subsection{Maintenance}

This document should be updated regularly to reflect:
\begin{itemize}
    \item New algorithm specifications
    \item Updated security recommendations
    \item New test vector sources
    \item Deprecated or broken algorithms
\end{itemize}

\vspace{2cm}

\begin{center}
\textbf{Last Updated:} \today\\
\textbf{Version:} 1.0\\
\textbf{Maintainer:} CryptoGL Documentation Team
\end{center}

\end{document} 

